\documentclass[12pt]{article}

% هوامش
\usepackage{geometry}
\geometry{margin=0.8in}

% ألوان/رؤوس/تذييل/تباعد
\usepackage{xcolor}
\usepackage{fancyhdr}
\usepackage{setspace}
\usepackage{enumitem}

% دعم يونيكود ولغات (يستلزم XeLaTeX)
\usepackage{fontspec}
\usepackage{polyglossia}

% Math packages for bmatrix/pmatrix/align etc.
\usepackage{amsmath, amssymb, mathtools}

% لصفّ الخيارات أفقياً مع لفّ تلقائي
\usepackage{tabularx, ragged2e, array}
\newcolumntype{Y}{>{\RaggedRight\arraybackslash}X}

% ماكرو جدول الاختيارات (A,B / C,D / E)
\newcommand{\renderchoices}[5]{%
  \begin{tabularx}{\linewidth}{@{}>{\bfseries}p{1.2em}Y >{\bfseries}p{1.2em}Y@{}}
    A.& #1 & B.& #2\\[0.25em]
    C.& #3 & D.& #4\\[0.25em]
    \multicolumn{1}{@{}>{\bfseries}p{1.2em}}{E.}& #5 & \multicolumn{2}{l}{}
  \end{tabularx}\par
}

\setmainlanguage[numerals=western]{arabic} % أرقام 0-9
\setotherlanguage{english}

% الخطوط
\newfontfamily\arabicfont[Script=Arabic,Scale=1.05]{Amiri}
\newfontfamily\englishfont{Times New Roman}

% ترويسة/تذييل
\pagestyle{fancy}
\fancyhf{}
\renewcommand{\headrulewidth}{0pt}
\cfoot{\textenglish{Page } \thepage}

% ماكرو الخيار (اختياري)
\newcommand{\choice}[2]{\textbf{#1}.~#2}

% فقرات
\setlength{\parskip}{0.4em}
\setlength{\parindent}{0pt}

\begin{document}
% ======= Header (fixed layout) =======
\noindent
\begin{tabular*}{\textwidth}{@{\extracolsep{\fill}} l c r}
\textenglish{Date: {\textenglish{September 4, 2025}}} &
{\Large \textbf{{\textenglish{Algorithms -1-}}}} &
{\textenglish{University of Aleppo}} \\
\textenglish{Student Name:} &
\textenglish{Model: {\textenglish{D}}} &
\textenglish{\textbf{Duration:} {\textenglish{90 minuts}}} \\
\textenglish{Student Number:} & & \\
\end{tabular*}

\vspace{0.9cm}

\section*{اختر الإجابة الصحيحة لكل سؤال}
\setstretch{1.25}

\noindent \textbf{\textenglish{1.}} {\textenglish{Test2}}

\renderchoices{\underline{{\textenglish{1}}}}{{\textenglish{2}}}{{\textenglish{3}}}{{\textenglish{4}}}{{\textenglish{5}}}
\hfill {\textenglish{[3 Points]}}\\[0.5em]

\noindent \textbf{\textenglish{2.}} {\textenglish{Solve for }}\(x \){\textenglish{ in the equation }}\(2x+6 = 0 \)

\renderchoices{{\textenglish{-1}}}{{\textenglish{-2}}}{\underline{{\textenglish{-3}}}}{{\textenglish{3}}}{}
\hfill {\textenglish{[1 Point]}}\\[0.5em]

\noindent \textbf{\textenglish{3.}} {\textarabic{احسب مشتق التابع }}\( f(x) = x^2\){\textarabic{ وأوجد نهايته عند اللانهاية }}\(\lim_{x \to \infty}f(x)\)

\renderchoices{{\textenglish{14}}}{{\textenglish{15}}}{\underline{{\textenglish{16}}}}{{\textenglish{17}}}{{\textenglish{17}}}
\hfill {\textenglish{[1 Point]}}\\[0.5em]

\noindent \textbf{\textenglish{4.}} {\textenglish{Hello!}}

\renderchoices{\underline{{\textenglish{1}}}}{{\textenglish{2}}}{{\textenglish{3}}}{{\textenglish{4}}}{}
\hfill {\textenglish{[2 Points]}}\\[0.5em]

\noindent \textbf{\textenglish{5.}} {\textenglish{Please calculate }}\(\int f(x)dx\)

\renderchoices{{\textenglish{13}}}{\underline{{\textenglish{14}}}}{{\textenglish{15}}}{{\textenglish{16}}}{}
\hfill {\textenglish{[2 Points]}}\\[0.5em]

\noindent \textbf{\textenglish{6.}} \(x=3+6 \){\textenglish{ what is the value of x?}}

\renderchoices{{\textenglish{7}}}{\underline{{\textenglish{9}}}}{{\textenglish{8}}}{{\textenglish{10}}}{{\textenglish{0}}}
\hfill {\textenglish{[1 Point]}}\\[0.5em]

\noindent \textbf{\textenglish{7.}} {\textenglish{Find the inverse of the following matrix: }}\(\begin{bmatrix} 1 & 3 & 4 \\ 5 & 2 & 1 \\ 6 & 8 & 1 \end{bmatrix}\)

\renderchoices{\underline{{\textenglish{There is no inverse for this matrix}}}}{\(\begin{pmatrix} 1 & 2 & 3 \\ 4 & 5 & 6 \\ 7 & 8 & 9  \end{pmatrix}\)}{{\textenglish{0}}}{{\textenglish{0}}}{}
\hfill {\textenglish{[4 Points]}}\\[0.5em]

\noindent \textbf{\textenglish{8.}} {\textenglish{Test3}}

\renderchoices{{\textenglish{1}}}{{\textenglish{2}}}{\underline{{\textenglish{3}}}}{{\textenglish{4}}}{{\textenglish{5}}}
\hfill {\textenglish{[1 Point]}}\\[0.5em]

\noindent \textbf{\textenglish{9.}} {\textenglish{Test}}

\renderchoices{{\textenglish{13}}}{\underline{{\textenglish{14}}}}{{\textenglish{15}}}{{\textenglish{16}}}{{\textenglish{17}}}
\hfill {\textenglish{[3 Points]}}\\[0.5em]

\noindent \textbf{\textenglish{10.}} {\textenglish{test4}}

\renderchoices{\underline{{\textenglish{1}}}}{{\textenglish{2}}}{{\textenglish{3}}}{{\textenglish{4}}}{{\textenglish{5}}}
\hfill {\textenglish{[1 Point]}}\\[0.5em]



\end{document}
